这里给出使用Schwinger-Dyson定理微扰展开的方法的展示, 我们不讨论具体的细节与一般性的推广. 因此我们只考虑有相互作用的实标量场, 我们有
\begin{theorem}[特化的Schwinger-Dyson定理]
    \begin{equation}
        (\Box_x+m^2)\braket{\phi_x\phi_1\cdots\phi_n}=\braket{\mathcal L_{\rm{int}}'(\phi_x)\phi_1\cdots\phi_n}-i\sum_j\delta_{xj}\braket{\phi_1\cdots\phi_{j-1}\phi_{j+1}\cdots\phi_n}
    \end{equation}
\end{theorem} 

然后我们展示通过这个定理计算$\phi^3$理论$\braket{\phi_1\phi_2}$. 
\begin{align}
    \braket{\phi_1\phi_2}&=\int\d^4x\delta_{1x}\braket{\phi_x\phi_2}=\int\d^4xi(\Box_x+m^2)D_{1x}\braket{\phi_x\phi_2}\\
    &=\int\d^4xiD_{1x}(\Box_x+m^2)\braket{\phi_x\phi_2}\\
    &=\int\d^4x\frac{ig}2 D_{1x}\braket{\phi_x^2\phi_2}-i\delta_{x2}\cdot iD_{1x}\\
    &=D_{12}+\frac{ig}2\int\d^4xD_{1x}\braket{\phi_x^2\phi_2}
\end{align}

而依葫芦画瓢我们有
\begin{align}
    \braket{\phi_x^2\phi_2}&=\int\d^4 iD_{y2}(\Box_y+m^2)\braket{\phi_x^2\phi_y}\\
    &=\frac{ig}2\int\d^4y D_{2y}\braket{\phi_x^2\phi_y^2}-2i\int\d^4y \delta_{xy}iD_{y2}\braket{\phi_x}\\
    &=\frac{ig}2\int\d^4y D_{2y}\braket{\phi_x^2\phi_y^2}+2D_{x2}\braket{\phi_x}
\end{align}

梅开三度
\begin{align}
    \braket{\phi_x}&=\int\d^4y iD_{yx}(\Box_y+m^2)\braket{\phi_y}\\
    &=\frac{ig}2\int\d^4y D_{yx}\braket{\phi_y^2}
\end{align}

于是最后整理得
\begin{equation}
    \braket{\phi_1\phi_2}=D_{12}-\left(\frac g2\right)^2\int \d^4x\d^4y(2D_{1x}D_{xy}D_{xy}D_{y2}+D_{1x}D_{xx}D_{yy}D_{y2}+2D_{1x}D_{xy}D_{yy}D_{x2})\label{2pt-correlation-expansion-schwinger-dyson}
\end{equation}

我们可以用Feynman图来表示这些二阶贡献, 并写下$\phi^3$理论在位形空间的Feynman规则.
\begin{figure}[htbp!]
    \begin{tikzpicture}
        \draw (-0.5,0) -- (1,0) ;
        \draw (2,0) -- (3.5,0);
        \filldraw[black] (1,0) circle (1pt) node[below] {$x$} ;
        \filldraw[black] (-0.5,0) circle (1pt)node[below] {$x_1$} ;
        \filldraw[black] (2,0) circle (1pt) node[below] {$y$} ;
        \filldraw[black] (3.5,0) circle (1pt)node[below] {$x_2$}  ;
        \draw (1,0)..controls(0.4,1.2)and(1.6,1.2)..(1,0);
        \draw (2,0)..controls(1.4,1.2)and(2.6,1.2)..(2,0);
        \draw (4.2,0)--(4.5,0);
        \draw (4.35,0.15)--(4.35,-0.15);
        \draw (5.2,0)--(6.7,0);
        \draw (7.7,0)--(9.2,0);
        \filldraw[black] (6.7,0) circle (1pt);% node[below] {$x$} ;
        \filldraw[black] (6.5,0)  node[below] {$x$} ;
        \filldraw[black] (5.2,0) circle (1pt)node[below] {$x_1$} ;
        \filldraw[black] (7.7,0) circle (1pt);
        \filldraw[black] (7.9,0) node[below] {$y$} ;
        \filldraw[black] (9.2,0) circle (1pt)node[below] {$x_2$}  ;
        \draw (6.7,0) arc (180:0:0.5 and 0.5);
        \draw (6.7,0) arc (180:360:0.5 and 0.5);
        \draw (9.9,0)--(10.2,0);
        \draw(10.05,0.15)--(10.05,-0.15);
        \draw (10.9,0)--(12.4,0);
        \draw (12.4,0)--(12.4,1);
        \filldraw[black] (12.4,0) circle (1pt) node[below] {$x$} ;
        \filldraw[black] (10.9,0) circle (1pt) node[below] {$x_1$} ;
        \filldraw[black] (12.4,1) circle (1pt) node[left] {$y$} ;
        \filldraw[black] (13.9,0) circle (1pt)node[below] {$x_2$}  ;
        \draw (12.4,0)--(13.9,0);
        \draw (12.4,1)..controls(13.6,1.6)and(13.6,0.4)..(12.4,1);
    \end{tikzpicture}
\end{figure}